\chapter{Lớp học và chiều sâu}
\markboth{Lớp học và chiều sâu}{Classroom and Depth}


\section{Đừng lấp đầy im lặng bằng lời vô nghĩa.}
\begin{truyen}{Đừng lấp đầy im lặng bằng lời vô nghĩa.}{Đừng lấp đầy im lặng bằng lời vô nghĩa.}
\chuhoa{C}{âu thầy nói câu ấy. Cả lớp im lặng — không ai nói, mỗi người đang nghĩ về điều vừa nghe.}
\textit{(The teacher said that. The class fell silent — no one spoke; everyone was thinking about what they had just heard.)}
\end{truyen}

\begin{baihoc}
Im lặng sau câu nói đúng lúc là dấu hiệu của tư duy.
\textit{(Silence after the right sentence is a sign of thinking.)}
\end{baihoc}

\begin{ghinhoanh}
\textbf{Short English to remember:}

Some sentences make the whole class think in silence.

\end{ghinhoanh}

\begin{tuhocnhanh}
1. Câu này khiến em nghĩ đến điều gì? \textit{What does this make you think of?}

2. Em có thể dùng câu này trong tình huống nào? \textit{When could you use this sentence?}
\end{tuhocnhanh}
\ngancach


\section{Suy nghĩ cần thời gian — và im lặng cho em thời gian đó.}
\begin{truyen}{Suy nghĩ cần thời gian — và im lặng cho em thời gian đó.}{Suy nghĩ cần thời gian — và im lặng cho em thời gi}
\chuhoa{C}{âu thầy nói câu ấy. Cả lớp im lặng — không ai nói, mỗi người đang nghĩ về điều vừa nghe.}
\textit{(The teacher said that. The class fell silent — no one spoke; everyone was thinking about what they had just heard.)}
\end{truyen}

\begin{baihoc}
Im lặng sau câu nói đúng lúc là dấu hiệu của tư duy.
\textit{(Silence after the right sentence is a sign of thinking.)}
\end{baihoc}

\begin{ghinhoanh}
\textbf{Short English to remember:}

Some sentences make the whole class think in silence.

\end{ghinhoanh}

\begin{tuhocnhanh}
1. Câu này khiến em nghĩ đến điều gì? \textit{What does this make you think of?}

2. Em có thể dùng câu này trong tình huống nào? \textit{When could you use this sentence?}
\end{tuhocnhanh}
\ngancach


\section{Người biết lắng nghe thường là người nói có trọng lượng sau này.}
\begin{truyen}{Người biết lắng nghe thường là người nói có trọng lượng sau ...}{Người biết lắng nghe thường là người nói có trọng }
\chuhoa{C}{âu thầy nói câu ấy. Cả lớp im lặng — không ai nói, mỗi người đang nghĩ về điều vừa nghe.}
\textit{(The teacher said that. The class fell silent — no one spoke; everyone was thinking about what they had just heard.)}
\end{truyen}

\begin{baihoc}
Im lặng sau câu nói đúng lúc là dấu hiệu của tư duy.
\textit{(Silence after the right sentence is a sign of thinking.)}
\end{baihoc}

\begin{ghinhoanh}
\textbf{Short English to remember:}

Some sentences make the whole class think in silence.

\end{ghinhoanh}

\begin{tuhocnhanh}
1. Câu này khiến em nghĩ đến điều gì? \textit{What does this make you think of?}

2. Em có thể dùng câu này trong tình huống nào? \textit{When could you use this sentence?}
\end{tuhocnhanh}
\ngancach


\section{Lớp học tốt không phải lớp ồn nhất hay im nhất — mà lớp biết khi nào nên im.}
\begin{truyen}{Lớp học tốt không phải lớp ồn nhất hay im nhất — mà lớp biết...}{Lớp học tốt không phải lớp ồn nhất hay im nhất — m}
\chuhoa{C}{âu thầy nói câu ấy. Cả lớp im lặng — không ai nói, mỗi người đang nghĩ về điều vừa nghe.}
\textit{(The teacher said that. The class fell silent — no one spoke; everyone was thinking about what they had just heard.)}
\end{truyen}

\begin{baihoc}
Im lặng sau câu nói đúng lúc là dấu hiệu của tư duy.
\textit{(Silence after the right sentence is a sign of thinking.)}
\end{baihoc}

\begin{ghinhoanh}
\textbf{Short English to remember:}

Some sentences make the whole class think in silence.

\end{ghinhoanh}

\begin{tuhocnhanh}
1. Câu này khiến em nghĩ đến điều gì? \textit{What does this make you think of?}

2. Em có thể dùng câu này trong tình huống nào? \textit{When could you use this sentence?}
\end{tuhocnhanh}
\ngancach


\section{Em không cần trả lời ngay; em cần nghĩ trước.}
\begin{truyen}{Em không cần trả lời ngay; em cần nghĩ trước.}{Em không cần trả lời ngay; em cần nghĩ trước.}
\chuhoa{C}{âu thầy nói câu ấy. Cả lớp im lặng — không ai nói, mỗi người đang nghĩ về điều vừa nghe.}
\textit{(The teacher said that. The class fell silent — no one spoke; everyone was thinking about what they had just heard.)}
\end{truyen}

\begin{baihoc}
Im lặng sau câu nói đúng lúc là dấu hiệu của tư duy.
\textit{(Silence after the right sentence is a sign of thinking.)}
\end{baihoc}

\begin{ghinhoanh}
\textbf{Short English to remember:}

Some sentences make the whole class think in silence.

\end{ghinhoanh}

\begin{tuhocnhanh}
1. Câu này khiến em nghĩ đến điều gì? \textit{What does this make you think of?}

2. Em có thể dùng câu này trong tình huống nào? \textit{When could you use this sentence?}
\end{tuhocnhanh}
\ngancach


\section{Im lặng là ngôn ngữ của suy nghĩ sâu.}
\begin{truyen}{Im lặng là ngôn ngữ của suy nghĩ sâu.}{Im lặng là ngôn ngữ của suy nghĩ sâu.}
\chuhoa{C}{âu thầy nói câu ấy. Cả lớp im lặng — không ai nói, mỗi người đang nghĩ về điều vừa nghe.}
\textit{(The teacher said that. The class fell silent — no one spoke; everyone was thinking about what they had just heard.)}
\end{truyen}

\begin{baihoc}
Im lặng sau câu nói đúng lúc là dấu hiệu của tư duy.
\textit{(Silence after the right sentence is a sign of thinking.)}
\end{baihoc}

\begin{ghinhoanh}
\textbf{Short English to remember:}

Some sentences make the whole class think in silence.

\end{ghinhoanh}

\begin{tuhocnhanh}
1. Câu này khiến em nghĩ đến điều gì? \textit{What does this make you think of?}

2. Em có thể dùng câu này trong tình huống nào? \textit{When could you use this sentence?}
\end{tuhocnhanh}
\ngancach


\section{Câu hỏi của thầy không phải để em trả lời nhanh — mà để em nghĩ kỹ.}
\begin{truyen}{Câu hỏi của thầy không phải để em trả lời nhanh — mà để em n...}{Câu hỏi của thầy không phải để em trả lời nhanh — }
\chuhoa{C}{âu thầy nói câu ấy. Cả lớp im lặng — không ai nói, mỗi người đang nghĩ về điều vừa nghe.}
\textit{(The teacher said that. The class fell silent — no one spoke; everyone was thinking about what they had just heard.)}
\end{truyen}

\begin{baihoc}
Im lặng sau câu nói đúng lúc là dấu hiệu của tư duy.
\textit{(Silence after the right sentence is a sign of thinking.)}
\end{baihoc}

\begin{ghinhoanh}
\textbf{Short English to remember:}

Some sentences make the whole class think in silence.

\end{ghinhoanh}

\begin{tuhocnhanh}
1. Câu này khiến em nghĩ đến điều gì? \textit{What does this make you think of?}

2. Em có thể dùng câu này trong tình huống nào? \textit{When could you use this sentence?}
\end{tuhocnhanh}
\ngancach


\section{Khi em im lặng, em đang tôn trọng chính suy nghĩ của mình.}
\begin{truyen}{Khi em im lặng, em đang tôn trọng chính suy nghĩ của mình.}{Khi em im lặng, em đang tôn trọng chính suy nghĩ c}
\chuhoa{C}{âu thầy nói câu ấy. Cả lớp im lặng — không ai nói, mỗi người đang nghĩ về điều vừa nghe.}
\textit{(The teacher said that. The class fell silent — no one spoke; everyone was thinking about what they had just heard.)}
\end{truyen}

\begin{baihoc}
Im lặng sau câu nói đúng lúc là dấu hiệu của tư duy.
\textit{(Silence after the right sentence is a sign of thinking.)}
\end{baihoc}

\begin{ghinhoanh}
\textbf{Short English to remember:}

Some sentences make the whole class think in silence.

\end{ghinhoanh}

\begin{tuhocnhanh}
1. Câu này khiến em nghĩ đến điều gì? \textit{What does this make you think of?}

2. Em có thể dùng câu này trong tình huống nào? \textit{When could you use this sentence?}
\end{tuhocnhanh}
\ngancach


\section{Nói trước khi nghĩ là đánh mất cơ hội hiểu.}
\begin{truyen}{Nói trước khi nghĩ là đánh mất cơ hội hiểu.}{Nói trước khi nghĩ là đánh mất cơ hội hiểu.}
\chuhoa{C}{âu thầy nói câu ấy. Cả lớp im lặng — không ai nói, mỗi người đang nghĩ về điều vừa nghe.}
\textit{(The teacher said that. The class fell silent — no one spoke; everyone was thinking about what they had just heard.)}
\end{truyen}

\begin{baihoc}
Im lặng sau câu nói đúng lúc là dấu hiệu của tư duy.
\textit{(Silence after the right sentence is a sign of thinking.)}
\end{baihoc}

\begin{ghinhoanh}
\textbf{Short English to remember:}

Some sentences make the whole class think in silence.

\end{ghinhoanh}

\begin{tuhocnhanh}
1. Câu này khiến em nghĩ đến điều gì? \textit{What does this make you think of?}

2. Em có thể dùng câu này trong tình huống nào? \textit{When could you use this sentence?}
\end{tuhocnhanh}
\ngancach


\section{Im lặng không phải đồng nghĩa với không biết — đôi khi là đang tìm câu trả lời đúng.}
\begin{truyen}{Im lặng không phải đồng nghĩa với không biết — đôi khi là đa...}{Im lặng không phải đồng nghĩa với không biết — đôi}
\chuhoa{C}{âu thầy nói câu ấy. Cả lớp im lặng — không ai nói, mỗi người đang nghĩ về điều vừa nghe.}
\textit{(The teacher said that. The class fell silent — no one spoke; everyone was thinking about what they had just heard.)}
\end{truyen}

\begin{baihoc}
Im lặng sau câu nói đúng lúc là dấu hiệu của tư duy.
\textit{(Silence after the right sentence is a sign of thinking.)}
\end{baihoc}

\begin{ghinhoanh}
\textbf{Short English to remember:}

Some sentences make the whole class think in silence.

\end{ghinhoanh}

\begin{tuhocnhanh}
1. Câu này khiến em nghĩ đến điều gì? \textit{What does this make you think of?}

2. Em có thể dùng câu này trong tình huống nào? \textit{When could you use this sentence?}
\end{tuhocnhanh}
\ngancach
